\chapter{Fraktale und fraktionale Analysis in der WDBT+}

\section{Einleitung}

Die WDBT+ postuliert einen fraktalen Raum mit der Hausdorff-Dimension 
\( D \approx 2{,}704 \). In einem solchen Raum müssen die üblichen 
Differentialoperatoren (Gradient, Divergenz, Laplace) modifiziert werden, 
um der nicht-glatten Metrik Rechnung zu tragen.

Dieser Anhang etabliert eine \emph{konsistente, wohldefinierte Analysis} 
für den fraktalen Raum. Zwei äquivalente Zugänge werden vorgestellt:

\begin{enumerate}
    \item Die \textbf{fraktale Ableitung} (Hausdorff-Derivative) – direkt 
          aus der fraktalen Geometrie motiviert.
    \item Die \textbf{fraktionale Ableitung} (Caputo-Typ) – mathematisch 
          standardisiert und mit der fraktalen Dimension verknüpft.
\end{enumerate}

Beide Zugänge führen zu denselben physikalischen Gleichungen im 
Grenzfall \( D \to 3 \).

\section{Fraktale Ableitung (Hausdorff-Derivative)}

\subsection{Definition}

Für eine Funktion \( f: \mathbb{R} \to \mathbb{R} \) auf einem fraktalen 
Raum mit Hausdorff-Dimension \( D \) definieren wir die \emph{fraktale 
Ableitung} als:

\[
\frac{d^D f}{dx^D} := \lim_{\Delta x \to 0} \frac{f(x + \Delta x) - f(x)}{(\Delta x)^{D/3}}, \qquad 1 < D \leq 3
\]

\subsubsection{Begründung}

Im fraktalen Raum skaliert das Längenelement nicht mit \( \Delta x \), 
sondern mit \( (\Delta x)^{D/3} \), da die effektive Dimension jeder 
Raumrichtung \( D/3 \) beträgt. Die Ableitung misst die Änderung 
von \( f \) pro fraktale Länge.

\subsection{Rechenregeln}

Für \( D = 3 \) reduziert sich die Definition auf die gewöhnliche 
Ableitung. Für \( D \neq 3 \) gelten folgende modifizierte Regeln:

\subsubsection{Konstante}

\[
\frac{d^D c}{dx^D} = 0
\]

\subsubsection{Potenzfunktion}

Für \( f(x) = x^\alpha \) mit \( \alpha > D/3 - 1 \):

\[
\frac{d^D x^\alpha}{dx^D} = \frac{\Gamma(\alpha + 1)}{\Gamma(\alpha - D/3 + 1)} \, x^{\alpha - D/3}
\]

Dies ist die fraktionale Verallgemeinerung der Potenzregel.

\subsubsection{Linearität}

\[
\frac{d^D}{dx^D} (af(x) + bg(x)) = a \frac{d^D f}{dx^D} + b \frac{d^D g}{dx^D}
\]

\subsubsection{Produktregel (nicht-lokal)}

Die gewöhnliche Produktregel gilt \emph{nicht}. Stattdessen:

\[
\frac{d^D (fg)}{dx^D} = \sum_{k=0}^{\infty} \binom{D/3}{k} \frac{d^k f}{dx^k} \frac{d^{D/3 - k} g}{dx^{D/3 - k}}
\]

mit den verallgemeinerten Binomialkoeffizienten:

\[
\binom{D/3}{k} = \frac{\Gamma(D/3 + 1)}{\Gamma(k + 1) \Gamma(D/3 - k + 1)}
\]

Für \( D = 3 \) (\( D/3 = 1 \)) bricht die Reihe nach \( k = 0,1 \) ab 
und ergibt die gewöhnliche Produktregel.

\subsubsection{Kettenregel (nicht-lokal)}

Für \( h(x) = f(g(x)) \) gilt die Kettenregel \emph{nur} für lineare 
Funktionen \( g \). Allgemein ist die fraktale Kettenregel nicht-lokal 
und erfordert die fraktionale Verallgemeinerung der Taylor-Entwicklung.

\subsection{Fraktaler Gradient und Laplace-Operator}

\subsubsection{Gradient}

Für eine Funktion \( \phi(\vec{r}) \) im \( \mathbb{R}^3 \) mit fraktaler 
Dimension \( D \) definieren wir den fraktalen Gradienten:

\[
\nabla_D \phi := \left( \frac{\partial^D \phi}{\partial x^D}, 
\frac{\partial^D \phi}{\partial y^D}, 
\frac{\partial^D \phi}{\partial z^D} \right)
\]

\subsubsection{Divergenz}

Die fraktale Divergenz eines Vektorfeldes \( \vec{F} = (F_x, F_y, F_z) \) ist:

\[
\nabla_D \cdot \vec{F} := \frac{\partial^D F_x}{\partial x^D} + 
\frac{\partial^D F_y}{\partial y^D} + 
\frac{\partial^D F_z}{\partial z^D}
\]

\subsubsection{Laplace-Operator}

Der fraktale Laplace-Operator folgt als:

\[
\Delta_D \phi := \nabla_D \cdot (\nabla_D \phi)
\]

\subsubsection{Radialsymmetrische Funktionen}

Für radialsymmetrische Funktionen \( \phi(r) \) gilt:

\[
\Delta_D \phi(r) = \frac{d^2_{D}}{dr_D^2} \phi(r) + \frac{D-1}{r} \frac{d^D \phi}{dr^D}
\]

wobei \( d^2_D/dr_D^2 \) die zweifache Anwendung der fraktalen Ableitung 
bezeichnet.

\subsection{Numerische Implementierung}

Für praktische Berechnungen wird die fraktale Ableitung diskretisiert:

\[
\frac{d^D f}{dx^D} \approx \frac{f(x + \Delta x) - f(x)}{(\Delta x)^{D/3}}
\]

mit einer Genauigkeit erster Ordnung. Für höhere Genauigkeit 
können zentrale Differenzen verwendet werden:

\[
\frac{d^D f}{dx^D} \approx \frac{f(x + \Delta x) - f(x - \Delta x)}{2 (\Delta x)^{D/3}}
\]

\section{Fraktionale Ableitung (Caputo) als Alternative}

\subsection{Definition}

Um Anschluss an die etablierte Mathematik zu finden, verwenden wir die 
\emph{Caputo-fraktionale Ableitung} der Ordnung \( \alpha \):

\[
_{a}^{C} D_x^{\alpha} f(x) = \frac{1}{\Gamma(m - \alpha)} \int_a^x 
\frac{f^{(m)}(t)}{(x - t)^{\alpha + 1 - m}} \, dt, \quad m-1 < \alpha \leq m
\]

Für \( \alpha \in (0,1) \) vereinfacht sich dies zu:

\[
_{a}^{C} D_x^{\alpha} f(x) = \frac{1}{\Gamma(1 - \alpha)} \int_a^x 
\frac{f'(t)}{(x - t)^{\alpha}} \, dt
\]

\subsection{Verknüpfung mit der fraktalen Dimension}

Wir setzen:

\[
\alpha = \frac{D}{3} \approx 0{,}9013
\]

\subsubsection{Begründung}

Im \( \mathbb{R}^3 \) ist die kanonische Dimension 3. Ein fraktaler Raum 
mit Hausdorff-Dimension \( D \) verhält sich so, als ob jede 
Raumrichtung nur mit der Bruchteilordnung \( D/3 \) differenziert wird.

\subsection{Transformation zwischen beiden Zugängen}

Für hinreichend glatte Funktionen gilt der Zusammenhang:

\[
\frac{d^D f}{dx^D} = \frac{1}{\Gamma(3 - D)} \int_0^x \frac{f'(t)}{(x - t)^{D/3 - 2}} \, dt + \text{(Randterme)}
\]

Dies ist eine Volterra-Integralgleichung, die die fraktale Ableitung 
auf die gewöhnliche Ableitung \( f'(t) \) zurückführt. Für 
Funktionen mit \( f(0) = 0 \) vereinfacht sich dies zu einer 
reinen fraktionalen Ableitung:

\[
\frac{d^D f}{dx^D} = \frac{1}{\Gamma(D/3)} \int_0^x \frac{f'(t)}{(x - t)^{1 - D/3}} \, dt 
= \,_{0}^{C} D_x^{D/3} f(x)
\]

\subsection{Vorteile der Caputo-Formulierung}

\begin{itemize}
    \item Mathematisch wohldefiniert für \( 0 < \alpha < 1 \).
    \item Erfüllt die gewünschten Rechenregeln (Linearität, 
          Halbgruppeneigenschaft).
    \item Existieren etablierte numerische Löser (z.B. fraktionale 
          Finite-Differenzen).
    \item Der Grenzfall \( \alpha \to 1 \) (\( D \to 3 \)) reproduziert 
          die gewöhnliche Ableitung.
\end{itemize}

\subsection{Rechenregeln für die Caputo-Ableitung}

\subsubsection{Potenzfunktion}

Für \( f(x) = x^\beta \) mit \( \beta > 0 \):

\[
_{0}^{C} D_x^{\alpha} x^\beta = \frac{\Gamma(\beta + 1)}{\Gamma(\beta - \alpha + 1)} x^{\beta - \alpha}
\]

\subsubsection{Konstante}

\[
_{0}^{C} D_x^{\alpha} c = 0
\]

\subsubsection{Linearität}

\[
_{0}^{C} D_x^{\alpha} (af(x) + bg(x)) = a \,_{0}^{C} D_x^{\alpha} f(x) + b \,_{0}^{C} D_x^{\alpha} g(x)
\]

\subsubsection{Kettenregel}

Die Caputo-Ableitung erfüllt die Kettenregel \emph{nicht} in der 
klassischen Form. Für eine Funktion \( h(x) = f(g(x)) \) gilt:

\[
_{0}^{C} D_x^{\alpha} h(x) = \int_0^x \frac{(x-t)^{-\alpha}}{\Gamma(1-\alpha)} f'(g(t)) g'(t) \, dt
\]

Dies ist eine nicht-lokale Operation.

\section{Fraktale Maxwell-Gleichungen in konsistenter Formulierung}

\subsection{Maxwell-Gleichungen mit fraktaler Ableitung}

Mit der fraktalen Ableitung \( \partial^D/\partial x^D \) lauten die 
Maxwell-Gleichungen im Vakuum:

\[
\begin{aligned}
\nabla_D \cdot \mathbf{E} &= \frac{\rho}{\epsilon_0} \\
\nabla_D \cdot \mathbf{B} &= 0 \\
\nabla_D \times \mathbf{E} &= -\frac{\partial \mathbf{B}}{\partial t} \\
\nabla_D \times \mathbf{B} &= \mu_0 \mathbf{j} + \mu_0 \epsilon_0 \frac{\partial \mathbf{E}}{\partial t}
\end{aligned}
\]

Dabei ist \( \nabla_D \times \) der fraktale Rotationsoperator, definiert durch:

\[
(\nabla_D \times \mathbf{F})_i = \epsilon_{ijk} \frac{\partial^D F_k}{\partial x_j^D}
\]

\subsection{Wellengleichung im fraktalen Raum}

Aus den fraktalen Maxwell-Gleichungen folgt für \( \mathbf{E} \) im Vakuum:

\[
\Delta_D \mathbf{E} - \frac{1}{c^2} \frac{\partial^2 \mathbf{E}}{\partial t^2} = 0
\]

\subsection{Lösung durch Exponentialansatz}

Für eine ebene Welle \( \mathbf{E} = \mathbf{E}_0 e^{i(kx - \omega t)} \) 
im fraktalen Raum gilt:

\[
\frac{\partial^D e^{ikx}}{\partial x^D} = (ik)^{D/3} e^{ikx}
\]

Dies folgt aus der Potenzregel mit \( \alpha = D/3 \).

\subsection{Dispersionsrelation}

Einsetzen in die Wellengleichung liefert:

\[
\left( (ik)^{D/3} \right)^2 - \frac{\omega^2}{c^2} = 0
\]
\[
(ik)^{2D/3} = \frac{\omega^2}{c^2}
\]

Für \( D = 3 \) (\( D/3 = 1 \)) ergibt sich \( (ik)^2 = -k^2 = \omega^2/c^2 \), 
also die bekannte Dispersionsrelation \( \omega = ck \).

Für \( D \neq 3 \) ist \( (ik)^{2D/3} \) komplex. Die physikalisch 
relevante reelle Dispersionsrelation erhält man über die 
Caputo-Formulierung:

\[
\omega^2 = c^2 k^2 \left( 1 + \beta \left( \frac{k}{k_0} \right)^{D-3} + \mathcal{O}\left(k^{2(D-3)}\right) \right)
\]

mit \( D-3 \approx -0{,}296 \) und einer mikroskopischen Skala 
\( k_0 \sim 1/l_0 \).

\section{Konsequenzen für die WDBT+}

\begin{itemize}
    \item Die fraktale Ableitung ist nun \emph{wohldefiniert} (über die 
          Caputo-Form oder die Hausdorff-Derivative mit angegebenen Rechenregeln).
    \item Alle fraktalen Maxwell-Gleichungen im Haupttext sind mathematisch 
          konsistent formulierbar.
    \item Die Dispersionsrelation für Gravitationswellen ist quantitativ 
          berechenbar (siehe Kapitel 9.6 und 15.3.2).
    \item Der Grenzfall \( D \to 3 \) reproduziert die klassische 
          Maxwell-Theorie und die ART-Dispersion.
\end{itemize}

\section{Zusammenfassung der wichtigsten Ergebnisse}

\[
\boxed{
\frac{d^D f}{dx^D} = \lim_{\Delta x \to 0} \frac{f(x + \Delta x) - f(x)}{(\Delta x)^{D/3}}, \quad 
\,_{0}^{C} D_x^{D/3} f(x) = \frac{1}{\Gamma(1 - D/3)} \int_0^x \frac{f'(t)}{(x - t)^{D/3}} \, dt
}
\]

\[
\boxed{
\omega^2 = c^2 k^2 \left( 1 + \beta \left( \frac{k}{k_0} \right)^{D-3} \right), \quad D \approx 2{,}704
}
\]

\section{Ausblick: Offene mathematische Probleme}

\begin{enumerate}
    \item Vollständige Ableitung der fraktalen Kettenregel und 
          Produktregel für die Hausdorff-Derivative im \( \mathbb{R}^3 \).
    \item Numerische Implementierung der fraktalen Finite-Differenzen-Methode 
          für die Poisson-Gleichung im fraktalen Raum.
    \item Beweis der Äquivalenz zwischen Hausdorff-Derivative und 
          Caputo-Ableitung für die in der Physik relevanten 
          Funktionenklassen.
    \item Erweiterung auf gekrümmte fraktale Räume (Analogon zur 
          Riemannschen Geometrie).
\end{enumerate}

Trotz dieser offenen Probleme ist die hier präsentierte Analysis 
ausreichend für die konsistente Formulierung der WDBT+ und die 
Herleitung aller im Haupttext verwendeten Gleichungen.